% GENERATED FILE. Edit canonical lessons, then run npm run generate:books.
% canonical-source-hash: fnv1a64:05cc9e9b3e546d48
% canonical-lessons: JA-C258-tanoshikatta, JA-C258-shourai, JA-C258-naritai, JA-C258-itta-koto, JA-C258-imamade, JA-W258-ima-kanji, JA-R258-how-it-felt

\chapter{How It Felt, and What I Hope For}
\label{ch:ja-how-it-felt-and-what-i-hope-for}
\hlchaptermodality{258}

\begin{chapteropening}
\textbf{By the end of this chapter:} \emph{I can say how something felt (tanoshikatta desu; omoshirokunakatta desu), what I want to become (isha ni naritai desu), what I have done before (itta koto ga arimasu), and that I never have until now (imamade ichido mo).}

Complete JA-R258-how-it-felt: say how a trip felt, what you hope to become and what you have done before.
\end{chapteropening}

\section[\texorpdfstring{\ja{たのしかった} … (tanoshikatta …)}{たのしかった … (tanoshikatta …)}]{\ja{たのしかった} / \ja{うれしかった} (tanoshikatta / ureshikatta) — how it felt}

\label{lesson:JA-C258-tanoshikatta}



\begin{quote}
Say \textquotedblleft{}fun\textquotedblright{} and \textquotedblleft{}glad\textquotedblright{}. (\emph{tanoshii}, \emph{ureshii}.) Then \textquotedblleft{}a trip\textquotedblright{}. (\emph{ryokō}.) Today they move into the past.
\end{quote}

\subsection*{You'll want to know: -\ja{かった}}
\begin{quote}
\textbf{\ja{りょこうは} \ja{たのしかったです。}}
\end{quote}

\begin{quote}
\emph{ryokō wa tanoshikatta desu.} — \textquotedblleft{}The trip was fun.\textquotedblright{}
\end{quote}

An i-adjective carries its own past: the last \textbf{\ja{い}} becomes \textbf{\ja{かった}}, and \textbf{\ja{です}} stays as it is to keep the sentence polite. \textbf{\ja{たのしい}} gives \textbf{\ja{たのしかった}} (\emph{tanoshikatta}), \textbf{\ja{うれしい}} gives \textbf{\ja{うれしかった}} (\emph{ureshikatta}).

\begin{quote}
\textbf{\ja{うれしかったです。}}
\end{quote}

\begin{quote}
\emph{ureshikatta desu.} — \textquotedblleft{}I was glad.\textquotedblright{}
\end{quote}

\begin{quote}
\textbf{\ja{きのうは} \ja{さむかったです。}}
\end{quote}

\begin{quote}
\emph{kinō wa samukatta desu.} — \textquotedblleft{}It was cold yesterday.\textquotedblright{}
\end{quote}

\textbf{\ja{いい}} takes its past from \textbf{\ja{よい}}, another form of the same word: \textbf{\ja{よかった}} (\emph{yokatta}), \textquotedblleft{}it was good\textquotedblright{}.

\begin{quote}
\textbf{\ja{えいがは} \ja{よかったです。}}
\end{quote}

\begin{quote}
\emph{eiga wa yokatta desu.} — \textquotedblleft{}The film was good.\textquotedblright{}
\end{quote}

\subsection*{You'll want to know: -\ja{くなかった}, and \ja{でした}}
\textquotedblleft{}Was not\textquotedblright{} comes from the \textbf{-\ja{くない}} you know: \textbf{\ja{おもしろくない}} becomes \textbf{\ja{おもしろくなかった}} (\emph{omoshirokunakatta}).

\begin{quote}
\textbf{\ja{えいがは} \ja{おもしろくなかったです。}}
\end{quote}

\begin{quote}
\emph{eiga wa omoshirokunakatta desu.} — \textquotedblleft{}The film was not interesting.\textquotedblright{}
\end{quote}

A na-adjective keeps the noun's way: \textbf{\ja{でした}} (\emph{deshita}) after it.

\begin{quote}
\textbf{\ja{うみは} \ja{しずかでした。}}
\end{quote}

\begin{quote}
\emph{umi wa shizuka deshita.} — \textquotedblleft{}The sea was calm.\textquotedblright{}
\end{quote}

\begin{quote}
\textbf{\ja{ちちは} \ja{げんきでした。}}
\end{quote}

\begin{quote}
\emph{chichi wa genki deshita.} — \textquotedblleft{}My father was well.\textquotedblright{}
\end{quote}

\subsection*{Your turn}
Say these aloud:

\begin{itemize}
\raggedright
  \item \textquotedblleft{}It was cold yesterday\textquotedblright{} — \emph{kinō wa samukatta desu}
  \item \textquotedblleft{}It was hot yesterday\textquotedblright{} — \emph{kinō wa atsukatta desu}
  \item \textquotedblleft{}It was cool yesterday\textquotedblright{} — \emph{kinō wa suzushikatta desu}
\end{itemize}

\begin{tcolorbox}[breakable,colback=teal!4,colframe=teal!35!black,title={Before you move on}]
\textquotedblleft{}It was sad\textquotedblright{}? (\textbf{\ja{かなしかったです}}.) \textquotedblleft{}It was good\textquotedblright{}? (\textbf{\ja{よかったです}}.) \textquotedblleft{}It was not fun\textquotedblright{}? (\textbf{\ja{たのしくなかったです}}.) \textquotedblleft{}It was quiet\textquotedblright{}? (\textbf{\ja{しずかでした}}.)
\end{tcolorbox}
\section[\texorpdfstring{\ja{しょうらい} (shōrai)}{しょうらい (shōrai)}]{\ja{しょうらい} (shōrai) — the future; one's prospects}

\label{lesson:JA-C258-shourai}



\begin{quote}
Before the new word: say \textquotedblleft{}next year\textquotedblright{} and \textquotedblleft{}worry\textquotedblright{}. (\emph{rainen}, \emph{shinpai}.)
\end{quote}

\subsection*{You'll want to know: \ja{しょうらい}}
\textbf{\ja{しょうらい}} — \emph{shōrai} — \textquotedblleft{}the future\textquotedblright{}, the years ahead of you: further off than \textbf{\ja{らいねん}}.

\begin{quote}
\textbf{\ja{しょうらいが} \ja{しんぱいです。}}
\end{quote}

\begin{quote}
\emph{shōrai ga shinpai desu.} — \textquotedblleft{}I worry about the future.\textquotedblright{}
\end{quote}

\begin{quote}
\textbf{\ja{しょうらい} \ja{りょこうに} \ja{いく} \ja{つもりです。}}
\end{quote}

\begin{quote}
\emph{shōrai ryokō ni iku tsumori desu.} — \textquotedblleft{}In the future I plan to travel.\textquotedblright{}
\end{quote}

\begin{cousinweb}
\textbf{\ja{しょうらい}} (\emph{shōrai}) is Sino-Japanese (kango): two Chinese characters read with their on'yomi, readings borrowed from Chinese. They are \emph{shō}, which KANJIDIC2 glosses \textquotedblleft{}from now on\textquotedblright{} among other senses, and \emph{rai}, \textquotedblleft{}come\textquotedblright{} (Middle Chinese \emph{loj}), the \emph{rai} of \textbf{\ja{らいねん}} (KANJIDIC2; Middle Chinese in Baxter's transcription). JMdict glosses the word \textquotedblleft{}future; (future) prospects\textquotedblright{} (JMdict).
\end{cousinweb}

\subsection*{Your turn}
Say these aloud:

\begin{itemize}
\raggedright
  \item \emph{shōrai}
  \item \textquotedblleft{}I worry about the future\textquotedblright{} — \emph{shōrai ga shinpai desu}
  \item \textquotedblleft{}In the future I plan to travel\textquotedblright{} — \emph{shōrai ryokō ni iku tsumori desu}
\end{itemize}

\begin{tcolorbox}[breakable,colback=teal!4,colframe=teal!35!black,title={Before you move on}]
\textquotedblleft{}The future\textquotedblright{}? (\textbf{\ja{しょうらい}}, \emph{shōrai}.) Nearer, \textquotedblleft{}next year\textquotedblright{}? (\textbf{\ja{らいねん}}.)
\end{tcolorbox}
\section[\texorpdfstring{\ja{〜}\ja{に} \ja{なりたいです} (ni naritai desu)}{〜に なりたいです (ni naritai desu)}]{\ja{〜}\ja{に} \ja{なりたいです} (ni naritai desu) — what I want to become}

\label{lesson:JA-C258-naritai}



\begin{quote}
Say \textquotedblleft{}my elder brother became a teacher\textquotedblright{}. (\emph{ani wa sensei ni narimashita}.) Then \textquotedblleft{}a doctor\textquotedblright{}. (\emph{isha}.)
\end{quote}

\subsection*{You'll want to know: \ja{〜}\ja{に} \ja{なりたいです}}
\begin{quote}
\textbf{\ja{いしゃに} \ja{なりたいです。}}
\end{quote}

\begin{quote}
\emph{isha ni naritai desu.} — \textquotedblleft{}I want to become a doctor.\textquotedblright{}
\end{quote}

\textbf{\ja{なる}} takes \textbf{-\ja{たい}}, as \textbf{\ja{いく}} does in \textbf{\ja{いきたい}}: \textbf{\ja{なりたい}} (\emph{naritai}), \textquotedblleft{}want to become\textquotedblright{}. \textbf{\ja{です}} makes it polite.

\begin{quote}
\textbf{\ja{しょうらい} \ja{せんせいに} \ja{なりたいです。}}
\end{quote}

\begin{quote}
\emph{shōrai sensei ni naritai desu.} — \textquotedblleft{}In the future I want to become a teacher.\textquotedblright{}
\end{quote}

A na-adjective takes \textbf{\ja{に}} too:

\begin{quote}
\textbf{\ja{日本語}\ja{が} \ja{じょうずに} \ja{なりたいです。}}
\end{quote}

\begin{quote}
\emph{nihongo ga jōzu ni naritai desu.} — \textquotedblleft{}I want to get good at Japanese.\textquotedblright{}
\end{quote}

\subsection*{You'll want to know: \ja{ゆめは} \ja{〜}\ja{ことです}}
\textbf{\ja{こと}} (\emph{koto}) is \textquotedblleft{}a thing, a matter\textquotedblright{}. After the dictionary form of a verb it turns the action into a thing, and \textbf{\ja{ゆめは}} (\emph{yume wa}), \textquotedblleft{}my dream is\textquotedblright{}, can name it:

\begin{quote}
\textbf{\ja{ゆめは} \ja{日本}\ja{に} \ja{すむ} \ja{ことです。}}
\end{quote}

\begin{quote}
\emph{yume wa nihon ni sumu koto desu.} — \textquotedblleft{}My dream is to live in Japan.\textquotedblright{}
\end{quote}

\begin{quote}
\textbf{\ja{ゆめは} \ja{うみの} \ja{ちかくに} \ja{すむ} \ja{ことです。}}
\end{quote}

\begin{quote}
\emph{yume wa umi no chikaku ni sumu koto desu.} — \textquotedblleft{}My dream is to live near the sea.\textquotedblright{}
\end{quote}

\subsection*{Your turn}
Say these aloud:

\begin{itemize}
\raggedright
  \item \textquotedblleft{}I want to become a teacher\textquotedblright{} — \emph{sensei ni naritai desu}
  \item \textquotedblleft{}I want to become a nurse\textquotedblright{} — \emph{kangoshi ni naritai desu}
  \item \textquotedblleft{}I want to become a singer\textquotedblright{} — \emph{kashu ni naritai desu}
\end{itemize}

\begin{tcolorbox}[breakable,colback=teal!4,colframe=teal!35!black,title={Before you move on}]
\textquotedblleft{}I want to become a doctor\textquotedblright{}? (\textbf{\ja{いしゃに} \ja{なりたいです}}.) \textquotedblleft{}My dream is to live in Japan\textquotedblright{}? (\textbf{\ja{ゆめは} \ja{日本}\ja{に} \ja{すむ} \ja{ことです}}.) What does \textbf{\ja{こと}} do after a verb? (It makes the action a thing.)
\end{tcolorbox}
\section[\texorpdfstring{\ja{いった} \ja{ことが} … (itta koto ga …)}{いった ことが … (itta koto ga …)}]{\ja{いった} \ja{ことが} \ja{あります} (itta koto ga arimasu) — I have done it before}

\label{lesson:JA-C258-itta-koto}



\begin{quote}
Say \textquotedblleft{}my dream is to live in Japan\textquotedblright{}. (\emph{yume wa nihon ni sumu koto desu}.) Then say \textquotedblleft{}once\textquotedblright{}. (\emph{ichido}.)
\end{quote}

\subsection*{You'll want to know: -\ja{た} \ja{ことが} \ja{あります}}
\begin{quote}
\textbf{\ja{日本}\ja{に} \ja{いった} \ja{ことが} \ja{あります。}}
\end{quote}

\begin{quote}
\emph{nihon ni itta koto ga arimasu.} — \textquotedblleft{}I have been to Japan.\textquotedblright{}
\end{quote}

The plain past, \textbf{\ja{こと}}, then \textbf{\ja{が} \ja{あります}} (\emph{ga arimasu}): \textquotedblleft{}there is the matter of having gone\textquotedblright{}. It says you have done something at least once, at some time in your life.

\begin{quote}
\textbf{\ja{すしを} \ja{たべた} \ja{ことが} \ja{あります。}}
\end{quote}

\begin{quote}
\emph{sushi o tabeta koto ga arimasu.} — \textquotedblleft{}I have eaten sushi.\textquotedblright{}
\end{quote}

\subsection*{You'll want to know: \ja{ありますか}, and \ja{いちども}}
To ask, \textbf{\ja{か}} goes at the end:

\begin{quote}
\textbf{\ja{うみで} \ja{およいだ} \ja{ことが} \ja{ありますか。}}
\end{quote}

\begin{quote}
\emph{umi de oyoida koto ga arimasu ka.} — \textquotedblleft{}Have you ever swum in the sea?\textquotedblright{}
\end{quote}

For \textquotedblleft{}never\textquotedblright{}, \textbf{\ja{いちども}} (\emph{ichido mo}), \textquotedblleft{}not even once\textquotedblright{}, and \textbf{\ja{ありません}}:

\begin{quote}
\textbf{\ja{いちども} \ja{およいだ} \ja{ことが} \ja{ありません。}}
\end{quote}

\begin{quote}
\emph{ichido mo oyoida koto ga arimasen.} — \textquotedblleft{}I have never swum.\textquotedblright{}
\end{quote}

\subsection*{Your turn}
Say these aloud:

\begin{itemize}
\raggedright
  \item \textquotedblleft{}I have eaten sushi\textquotedblright{} — \emph{sushi o tabeta koto ga arimasu}
  \item \textquotedblleft{}I have eaten tempura\textquotedblright{} — \emph{tenpura o tabeta koto ga arimasu}
  \item \textquotedblleft{}I have eaten natto\textquotedblright{} — \emph{nattō o tabeta koto ga arimasu}
\end{itemize}

\begin{tcolorbox}[breakable,colback=teal!4,colframe=teal!35!black,title={Before you move on}]
\textquotedblleft{}I have been there\textquotedblright{}? (\textbf{\ja{いった} \ja{ことが} \ja{あります}}.) \textquotedblleft{}Have you ever swum?\textquotedblright{} (\textbf{\ja{およいだ} \ja{ことが} \ja{ありますか}}.) \textquotedblleft{}I have never swum\textquotedblright{}? (\textbf{\ja{いちども} \ja{およいだ} \ja{ことが} \ja{ありません}}.)
\end{tcolorbox}
\section[\texorpdfstring{\ja{今}\ja{まで} (imamade)}{今まで (imamade)}]{\ja{今}\ja{まで} (imamade) — until now}

\label{lesson:JA-C258-imamade}



\begin{quote}
Before the new word: say \textquotedblleft{}now\textquotedblright{}. (\emph{ima}.) Then \textquotedblleft{}I have never swum\textquotedblright{}. (\emph{ichido mo oyoida koto ga arimasen}.)
\end{quote}

\subsection*{You'll want to know: \ja{今}\ja{まで}}
\textbf{\ja{今}\ja{まで}} — \emph{imamade} — \textquotedblleft{}until now\textquotedblright{}, \textquotedblleft{}so far\textquotedblright{}: \textbf{\ja{いま}}, \textquotedblleft{}now\textquotedblright{}, with \textbf{\ja{まで}}, \textquotedblleft{}as far as\textquotedblright{}.

\begin{quote}
\textbf{\ja{今}\ja{まで} \ja{いちども} \ja{いった} \ja{ことが} \ja{ありません。}}
\end{quote}

\begin{quote}
\emph{imamade ichido mo itta koto ga arimasen.} — \textquotedblleft{}I have never been there until now.\textquotedblright{}
\end{quote}

\begin{quote}
\textbf{\ja{今}\ja{まで} \ja{ちかくに} \ja{すんでいました。}}
\end{quote}

\begin{quote}
\emph{imamade chikaku ni sunde imashita.} — \textquotedblleft{}Until now I lived nearby.\textquotedblright{}
\end{quote}

The first sign is new: the kanji for \textbf{\ja{いま}}, read \textbf{\ja{いま}} (\emph{ima}) here. The next lesson writes it.

\begin{cousinweb}
\textbf{\ja{今}\ja{まで}} (\emph{imamade}) is a native Japanese word (wago), inherited rather than borrowed from Chinese: \textbf{\ja{いま}} (\emph{ima}), \textquotedblleft{}now\textquotedblright{}, written with its kanji and read the native way, and \textbf{\ja{まで}} (\emph{made}), \textquotedblleft{}as far as\textquotedblright{}. JMdict glosses it \textquotedblleft{}until now; so far\textquotedblright{}, and KANJIDIC2 gives \ja{いま} as the kanji's native reading (JMdict, KANJIDIC2).
\end{cousinweb}

\subsection*{Your turn}
Say these aloud:

\begin{itemize}
\raggedright
  \item \emph{imamade}
  \item \textquotedblleft{}I have never been there until now\textquotedblright{} — \emph{imamade ichido mo itta koto ga arimasen}
  \item \textquotedblleft{}I have never eaten it until now\textquotedblright{} — \emph{imamade ichido mo tabeta koto ga arimasen}
\end{itemize}

\begin{tcolorbox}[breakable,colback=teal!4,colframe=teal!35!black,title={Before you move on}]
\textquotedblleft{}Until now\textquotedblright{}? (\textbf{\ja{今}\ja{まで}}, \emph{imamade}.) How is its first sign read? (\textbf{\ja{いま}}, \emph{ima}.)
\end{tcolorbox}
\section[\texorpdfstring{\ja{今} (ima)}{今 (ima)}]{\ja{今} — now, four strokes}

\label{lesson:JA-W258-ima-kanji}



\begin{quote}
Say \emph{imamade}, until now. Then call up \textbf{\ja{人}}: two strokes, the left one first.
\end{quote}

\subsection*{Script — the shape}
\begin{quote}
\textbf{\ja{今}} — \emph{ima}
\end{quote}

Its parts:

\begin{itemize}
\raggedright
  \item \ja{人} (a roof) on top
  \item a short bar
  \item a bar turning down to the left
\end{itemize}

Four strokes, three lifts:

\begin{enumerate}
\raggedright
  \item draw the falling stroke on the left of the roof, from the top down to the left
  \item lift, draw the falling stroke on the right of the roof, down to the right
  \item lift, draw the short bar under the roof from left to right
  \item lift, draw the lower bar from left to right, then without lifting turn down to the left
\end{enumerate}

The roof is \textbf{\ja{人}}, in its order, left then right, but here the second stroke starts at the very top and both spread wide. Here it is read \textbf{\ja{いま}} (\emph{ima}), as in \textbf{\ja{今}\ja{まで}}; this lesson asks only for that reading.

\subsection*{Your turn}
\begin{enumerate}
\raggedright
  \item Find \textbf{\ja{今}} in \textbf{\ja{今}\ja{まで}} and say \emph{imamade}.
  \item Copy \textbf{\ja{今}} three times, saying \emph{ima} at each one.
  \item Find it once more, in \textbf{\ja{今}\ja{まで} \ja{ちかくに} \ja{すんでいました}} (\emph{imamade chikaku ni sunde imashita}), and say the line.
\end{enumerate}

\begin{tcolorbox}[breakable,colback=teal!4,colframe=teal!35!black,title={Before you move on}]
Which sign is its roof? (\textbf{\ja{人}}, left stroke first.)

Source: \href{https://github.com/KanjiVG/kanjivg/blob/master/kanji/04eca.svg}{KanjiVG stroke-order data}: KanjiVG, kanji/04eca.svg, kanji \ja{今} (U+4ECA), 4 stroke paths kvg:04eca-s1 to s4 (Ulrich Apel and contributors, CC BY-SA 3.0). Only the order and direction are taken from it.
\end{tcolorbox}
\section[(review)]{Review — how it felt, and what I hope for}

\label{lesson:JA-R258-how-it-felt}



\begin{quote}
Say \textquotedblleft{}the trip was fun\textquotedblright{}. (\emph{ryokō wa tanoshikatta desu}.) Then \textquotedblleft{}I want to become a teacher\textquotedblright{}. (\emph{sensei ni naritai desu}.)
\end{quote}

\subsection*{Your turn}
Say these aloud:

\begin{itemize}
\raggedright
  \item It was cold yesterday.
  \item The film was not interesting.
  \item In the future I want to become a doctor.
  \item My dream is to live near the sea.
  \item Have you ever eaten sushi?
\end{itemize}

\begin{tcolorbox}[breakable,colback=teal!4,colframe=teal!35!black,title={Before you move on}]
\textquotedblleft{}I was glad\textquotedblright{}? (\textbf{\ja{うれしかったです}}.) \textquotedblleft{}It was good\textquotedblright{}? (\textbf{\ja{よかったです}}.) \textquotedblleft{}I want to become\textquotedblright{}? (\textbf{\ja{〜}\ja{に} \ja{なりたいです}}.) \textquotedblleft{}I have done it\textquotedblright{}? (\textbf{-\ja{た} \ja{ことが} \ja{あります}}.) \textquotedblleft{}Until now\textquotedblright{}? (\textbf{\ja{今}\ja{まで}}.)
\end{tcolorbox}
